\documentclass[10pt,a4paper]{amsart}
\usepackage[margin=19mm]{geometry}
\usepackage{amsmath,amssymb,mathtools}
\usepackage[scr=rsfs,cal=euler]{mathalfa}
\usepackage{xfrac}
\usepackage[unicode,hidelinks,bookmarks=false]{hyperref}
\newcommand{\ie}{i.e.\ }
\newcommand{\cf}{cf.\ }
\newcommand{\resp}{respectively\ }
\newcommand{\Subsection}{\subsection}
\providecommand{\nobreakdash}{\nobreak\hbox{-}}
\setlength{\emergencystretch}{2em}
\multlinegap0pt
\usepackage{bm}
\usepackage{xcolor}
\usepackage{amssymb}
\usepackage{dsfont}
\usepackage{mathtools}
\usepackage{float}
\usepackage{tikz}
\usepackage{enumerate}
\newtheorem{lemma}{Lemma}
\newcommand{\N}{\mathbb{N}}
\newcommand{\NN}{\mathbb{N}}
\newcommand{\R}{\mathbb{R}}
\newcommand{\RR}{\mathbb{R}}
\newcommand{\Z}{\mathbb{Z}}
\renewcommand{\hat}{\widehat}
\DeclareMathOperator{\meas}{meas}
\newcommand{\ua}{\vec{\mathrm{a}}}
\newcommand{\ur}{\vec{\mathrm{r}}}
\newcommand{\uv}{{\vec{\mathrm{v}}}}
\newcommand{\uw}{{\vec{\mathrm{w}}}}
\newcommand{\ux}{{\vec{\mathrm{x}}}}
\title{A two-dimensional delta symbol method and its application to pairs of quadratic forms}
\author{Junxian Li, Simon L. Rydin Myerson, Pankaj Vishe}
\date{Source: arXiv:2411.11355v2 (CC BY 4.0)}
\begin{document}
\maketitle

\noindent\emph{Source and licence.} A two-dimensional delta symbol method and its application to pairs of quadratic forms. Version arXiv:2411.11355v2 (CC BY 4.0), distributed by arXiv under the Creative Commons Attribution 4.0 International licence, http://creativecommons.org/licenses/by/4.0/, as recorded in the arXiv licence metadata for this exact version and shown on the abstract page https://arxiv.org/abs/2411.11355v2. Full licence notice: \texttt{licenses/arxiv-2411.11355-CC-BY-4.0.txt}.\par\smallskip

\noindent\textit{Editorial scope.} Counted unit: the article's lemma lem:1 (source0.tex lines 1711--1723). The article's complete original proof of it is delivered with it (source0.tex lines 1725--1898, one \texttt{proof} environment of 174 lines). No mathematics is written, repaired or re-derived: the statement and the proof are the article's own lines, in the article's own notation, and no retained line is substituted or re-flowed.  Retained uncounted as the source-local context the proof needs: lines 809--815; lines 822--837; lines 1020--1045; lines 1492--1507; lines 1511--1517; lines 1520--1527; lines 1558--1585.  Nothing was omitted from any retained span, so the package carries no omission record.  No retained line contains a citation, so the wrapper carries no bibliography and no bibliography entry.  No retained line contains a figure environment, an image-inclusion command or drawing code, so no image is delivered and the package declares no asset dependency.  Editorial formatting only, in addition: an AMS \texttt{amsart} preamble that restates the article's own declarations and macros used by the retained lines, namely the environments \texttt{lemma} on separate counters, together with the standard packages those lines need.  The retained environments print in this document as Lemma 1 in order of appearance, whereas the article prints the same items with the numbering of its own full document.  The complete native source of the article as distributed in the arXiv e-print for this exact version is bundled unchanged as \texttt{sources/168661/source0.tex}.
\begin{equation}
	\label{eq:hyzdef}
	h_{\omega_1}(y,z)
	=
	\sum_{j\in\N}
	\frac{1}{yj}\left(\omega_1(yj)-\omega_1\left(\frac{|z|}{yj}\right)\right).
\end{equation}

\begin{lemma}\label{lem:HB}
	Let $\omega_1$ be a
    smooth function supported on $(1/2, 1)$. Recall the function $h_{\omega_1}(y,z)$ from \eqref{eq:hyzdef}.
    \begin{enumerate}
        \item 
    The value of $h_{\omega_1}(y,z)$ is non-zero only if  $y\leq \max (1, 2|z|)$. Moreover $h_{\omega_1}(y,z)\ll_{\omega_1} 1/y$ for all $z$.
    \item
    Suppose that \(A,B, \delta>0\) and that \(f:\R\to \R\) satisfies \(|f^{(k)}(z)|\leq C_k A^k\) for all \(z\in \R\) and some sequence of positive real numbers \(\boldsymbol{C}=(C_k)_{k\geq 0}\). 
    Then for \(0<y\leq Q^{-\delta}\min\{B/A,1\}\) we have
	\[
	\int_{\R} f(z)h_{\omega_1}(y,Bz)\,dz
	= B^{-1}f(0)\int_{\R}\omega_1(z)\,dz
	+O_{N,\delta,\boldsymbol{C},\omega_1}(B^{-1}Q^{-N}).
	\]
    \end{enumerate}
\end{lemma}

		\begin{align}
		\label{defp2}\quad
			p_{1,q}(\uw)
			&=
			-\frac{2c}{Q}
			\int_{\R^2} 
			\omega_0\left(\frac{|\ux|}{Q^{3/2}}\right)
			\sum_{j\in\NN}
			\frac{1}{q^2j^2}
			\omega
			\left(\frac{|\ux|}{jqQ^{1/2}}
			\right)
			e(-\uw\cdot \ux)
			\,d\ux,\\
				\label{defp1}\quad
				p_{2,\ur,k,q}(\uw)
				&=\frac{c}{Q^3}
				\int_{\R^2} 
				\omega_0\left(\frac{|\ux|}{Q^{3/2}}\right)
				h\left(
				\frac{kq}{Q},
				\frac{\ur\cdot \ux }{ Q^2}
				\right)
				e(-\uw\cdot \ux)
				\,d\ux.
	\end{align}

\begin{lemma}\label{lem:minkowski_basis}
	Let $\Lambda$ be a lattice in $\RR^2$.
	Let \(M\) be an \(m \times 2\) real matrix with full rank. Further let 
	\( \ux_1 \in \Lambda \) with \(|M\ux_1| = \mu_M(\Lambda )=\min\{ |M\ux| : \ux \in \Lambda, \ux \neq\vec{0}\}\). Then there is an \(\ux_2\in \Lambda \) satisfying the following
	\begin{align*}\
	\Lambda & = \Z \ux_1 + \Z \ux_2,
	\\
	| r_1 M\ux_1 + r_2 M\ux_2 |
	&\asymp_m |r_1||M\ux_1| + |r_2| |M\ux_2|,
	\quad (\ur\in\R^2),
	\\
	|  M\ux_1 || M\ux_2 |
	&\asymp_m
	\mathrm{covol}(\Lambda)\meas\{ M [0,1]^2 \}.
	\end{align*}
\end{lemma}

	\begin{equation}\label{def:Mmatrdef_0}
		M\ur
		=
		\left(\frac{r_1}{Q^{1/2}},
		\frac{r_2}{Q^{1/2}},
		Q\uw\cdot \ur^\perp\right)^T.
	\end{equation}

\begin{lemma}
	\label{lem:matrix_M}
	 For $M$ in defined in \eqref{def:Mmatrdef_0}, we have 
	\begin{equation*}
		%\label{eq:applying_minkowski_basis2}
		\meas \big\{ M [0,1]^2 \big\} \asymp  Q^{1/2}|\uw|+Q^{-1}.
	\end{equation*}
\end{lemma} 

\begin{lemma}\label{lem:p1_sum_asymptotic}
	Let $\delta>0$ and $M=M(\uw)$ be as in \eqref{def:Mmatrdef_0} and let $\mu_{M}=\mu_M(\Lambda(\ua, q))$ the norm of the smallest non-zero vector in $M (\Lambda (\ua, q))$. If
	\begin{equation}\label{eq:muM assump}
	\mu_M\geq
	Q^{\delta}(qQ^{-1}+|\uw|qQ^{1/2}),
	\end{equation}
	then
	\begin{equation*}%\label{eq:p1_sum_asymptotic}
    \begin{split}
	&\sum_{{\ur\in\Lambda(\ua,q)}}
	\omega\left(
	\frac{ |\ur| }{ Q^{1/2}}
	\right)
	p_{2,\ur,k,q}(\uw)+O_{N,\delta}(Q^{-N})
	\\=&
	\frac{cQ^{1/2}}{q}
	\int_{\R^2}\frac{1}{|\ur|}
	\omega\left(
	\frac{ |\ur| }{ Q^{1/2}}
	\right)
	\hat{\omega}_0
	\Big(
	\frac{Q^{3/2}}{|\ur|} \uw\cdot\ur^\perp
	\Big)\,d\ur
    \end{split}
	\end{equation*}
	for any $N>0$.
\end{lemma}

\begin{lemma}\label{lem:1}
Let $M=M(\uw)$ and $\mu_{M}=\mu_M (\Lambda(\ua, q))$ be as in Lemma~\ref{lem:p1_sum_asymptotic}. Let $1\leq q\leq Q$ and  $\delta>0$. If
	\[
	\mu_M\geq
	Q^{\delta}(qQ^{-1}+|\uw|qQ^{1/2}),
	\]
	then
	\[
	p_{\Lambda(\ua,q)}(\uw)
	=1+O_{N,\delta}(Q^{-N}),
	\]
	for any $N>0$.
\end{lemma}

\begin{proof}
By Lemma~\ref{lem:p1_sum_asymptotic}, and making a change of variables using \(\ur=z'\uv_{\theta}, \ \uv_{\theta}=(\cos \theta, \sin \theta),\  z'=|\ur|\), we obtain
	\begin{align*}
		&\quad p_{\Lambda(\ua, q)}(\uw)-
		p_{1,q}(\uw)\\&=\frac{cQ^{1/2}}{q}
	\int_{\R^2}\frac{1}{|\ur|}
	\omega\left(
	\frac{ |\ur| }{ Q^{1/2}}
	\right)
	\hat{\omega}_0
	\Big(
	\frac{Q^{3/2}}{|\ur|} \uw\cdot\ur^\perp
	\Big)\,d\ur
	+O_{N,\delta}(Q^{-N})\\
    &=\frac{cQ^{1/2}}{q}\int_{0}^\infty \omega\Big(\frac{z'}{Q^{1/2}}\Big) dz' \int_{0}^{2\pi}\hat{\omega}_0\Big(
     Q^{3/2}\uw \cdot \uv_{\theta}^\perp \Big)d\theta+O_{N, \delta}(Q^{-N})\\
     &=\frac{cQ}{q}\int_{0}^{2\pi}\int_{\RR}\omega_0(z)e\Big(-zQ^{3/2}\uw 
     \cdot \uv_{\theta}^\perp\Big)dz d\theta +O_{N, \delta}(Q^{-N}).
    \end{align*}
\iffalse
\textcolor{gray}{	By Lemma~\ref{lem:p1_sum_asymptotic}, and making a change of variables using \(\ux'=z\ur/|\ur|, z'=|\ur|\), we obtain
 \begin{align*}
		p_{\Lambda(\ua, q)}(\uw)-
		p_{1,q}(\uw)&=\frac{cQ^{1/2}}{q}
	\int_{\R^2}\frac{1}{|\ur|}
	\omega\left(
	\frac{ |\ur| }{ Q^{1/2}}
	\right)
	\hat{\omega}_0
	\Big(
	\frac{Q^{3/2}}{|\ur|} \uw\cdot\ur^\perp
	\Big)\,d\ur
	+O_{N,\delta}(Q^{-N})\\
		&=
		\frac{cQ^{1/2}}{q}
		\int_{\R^2}
		\int_{\R}
		\frac{1}{|\ur|}
		\omega\left(
		\frac{ |\ur| }{ Q^{1/2}}
		\right)
		{\omega}_0(z)
		e\left(-zQ^{3/2}\uw\cdot\frac{ \ur^\perp}{|\ur|}\right)
		\,dz
		\,d\ur+O_{N,\delta}(Q^{-N}).
	\end{align*}
	for any $N>0$. 
}   
\fi
By Lemmas~\ref{lem:minkowski_basis} and~\ref{lem:matrix_M}, we have 
\begin{align}\label{muM}
\mu_M \ll (qQ^{-1}+|\uw|qQ^{1/2})^{1/2}.
\end{align} Therefore, by the assumption on \(\mu_M\) we have
	 \(q\ll Q^{1-\delta}\), and therefore upon substituting  \(z=Q^{3/2}z\) we get 
	\begin{align*}
		&\quad p_{\Lambda(\ua, q)}(\uw)-
		p_{1,q}(\uw)+O_{N,\delta}(Q^{-N})
		\\&=
		\frac{c}{Q^{3/2}}
		\left(\sum_{j\in\N} \omega\left(\frac{qj}{Q}\right)\right)
		\int_{0}^{2\pi}\int_{\RR}
		{\omega}_0\left(\frac{z}{Q^{3/2}}\right)
		e\left(-zQ^{3/2}\uw\cdot \uv_\theta^\perp\right)dzd\theta
		.
	\end{align*}
    
	Note that by making a change of variable $\ux \mapsto \ux^\perp$ in $p_{1, q}(\uw)$ defined in \eqref{defp2}, and further changing to polar coordinates, we can write
    \begin{equation}
    \begin{split}
    %\label{defp2repeated}
    &p_{1,q}(\uw)
			=
			-\frac{2c}{Q}
			\int_{\R^2} 
			\omega_0\left(\frac{|\ux|}{Q^{3/2}}\right)
			\sum_{j\in\NN}
			\frac{1}{q^2j^2}
			\omega
			\left(\frac{|\ux|}{jqQ^{1/2}}
			\right)
			e(-\uw\cdot \ux^\perp)
			\,d\ux\\
            &=-\frac{2c}{Q}\int_{0}^{2\pi}\int_{0}^\infty z\omega_0\Big(\frac{z}{Q^{3/2}}\Big) \sum_{j\in\NN}
			\frac{1}{q^2j^2}
			\omega
			\left(\frac{z}{jqQ^{1/2}}
			\right)
			e\Big(-z\uw\cdot \uv_{\theta}^\perp\Big)
			\,dz d\theta.\label{eq:p1qmanipulation}
            \end{split}
            \end{equation}

It is straight-forward to see, by replacing $\theta$ with $\theta+\pi$, that
\begin{align*}
    &\int_{0}^{2\pi}\int_{0}^\infty z\omega_0\Big(\frac{z}{Q^{3/2}}\Big) \sum_{j\in\NN}
			\frac{1}{q^2j^2}
			\omega
			\left(\frac{z}{jqQ^{1/2}}
			\right)
			e\Big(-z\uw\cdot \uv_{\theta}^\perp\Big)
			\,dz d\theta\\&=\int_{0}^{2\pi}\int_{0}^\infty z\omega_0\Big(\frac{z}{Q^{3/2}}\Big) \sum_{j\in\NN}
			\frac{1}{q^2j^2}
			\omega
			\left(\frac{z}{jqQ^{1/2}}
			\right)
			e\Big(z\uw\cdot \uv_{\theta}^\perp\Big)
			\,dz d\theta.
\end{align*}
As $\omega_0$ is even, we may therefore re-write \eqref{eq:p1qmanipulation} as
\begin{align*}p_{1,q}(\uw)
			=-\frac{c}{Q}\int_{0}^{2\pi}\int_{\mathbb R} |z|\omega_0\Big(\frac{z}{Q^{3/2}}\Big) \sum_{j\in\NN}
			\frac{1}{q^2j^2}
			\omega
			\left(\frac{|z|}{jqQ^{1/2}}
			\right)
			e\big(-z\uw\cdot \uv_{\theta}^\perp\big)
			\,dz d\theta.
            \end{align*}

 Therefore, we have 
	\begin{align*}
		&\quad \quad p_{\Lambda(\ua, q)}(\uw)+O_{N, \delta}(Q^{-N})\\&=\frac{c}{Q^{\frac{3}{2}}}
		\int_{0}^{2\pi}\int_{\RR}\Big(\sum_{j\in\N} \omega\big(\frac{qj}{Q}\big)-\frac{Q^{\frac{1}{2}}|z|}{j^2q^2} \omega\big(\frac{|z|}{jqQ^{1/2}}\big)\Big)
		\omega_0\big(\frac{z}{Q^{3/2}}\big)
		e(-z\uw\cdot \uv_{\theta}^\perp)
		dzd\theta,	
	\\	&=
		\frac{c}{Q^{\frac{3}{2}}}
		\int_{0}^{2\pi}\int_{\RR} 
	\omega_0\big(\frac{z}{Q^{3/2}}\big) 
		h_2\big(
		\frac{q}{Q}
		,
		\frac{z}{Q^{3/2}}
		\big)
		e\big(-z\uw\cdot \uv_{\theta}^\perp\big)
		dzd\theta,
	\end{align*}
	where
	\begin{align*}
		h_2(y,z)
		&=
		\sum_{j\in\N}
		\frac{1}{yj}
		\left(
		{yj}
		\omega\left(yj\right)
		-
		\frac{|z|}{yj}
		\omega\left(\frac{|z|}{yj}\right)
		\right).
	\end{align*}
	Using the assumption in the lemma and \eqref{muM},  we have 
	\(qQ^{-1}+|\uw|qQ^{1/2} \ll Q^{-2\delta}\), so that we can apply Lemma~\ref{lem:HB} to the $z$-integral with
	\begin{align*}
		\omega_1(x)={x\omega(x)},\quad
		y=\frac{q}{Q},\quad
		A=Q^{-3/2}+|\uw|,\quad
		\\B=Q^{-3/2},\quad
		f(z)=\omega_0\Big(\frac{z}{Q^{3/2}}\Big)e\Big(-z\uw\cdot \uv_{\theta}^\perp\Big),
	\end{align*}
	and conclude that
	\begin{equation*}\begin{split}%\label{eq:22}
		p_{\Lambda(\ua, q)}(\uw)
		&=c \cdot \omega_0(0)e(0)
		\int_{0}^{2\pi} 
		\int_0^\infty x
		\omega(x)\,dx
		\,d\theta+O_{N,\delta}(Q^{-N})
		\\&=1+O_{N,\delta}(Q^{-N}),
        \end{split}
	\end{equation*}
 which proves the lemma.
\end{proof}

\end{document}
